\section{Consume only the keystream bytes the message needs}
Once the two block functions agree, the outer transformation has a smaller task.
Copy the starting state into private working storage. Generate a block from that
copy, then XOR at most 64 message bytes with its 64 keystream bytes. A final message
fragment may use only one byte. Do not load the other 63 source bytes merely because
a whole keystream block was generated.

Keep source, destination and remaining count in saved registers across the block
call. Choose \(m=\min(64,\text{remaining})\). For each byte, load both operands
before storing the result. This ordering also supports exact in-place operation:
the store changes only the byte whose original input has already been read.
\Needspace{9\baselineskip}
\MARSCodeLabel{Transform one byte; advance both pointers afterward.}
\begin{minted}[samepage=true]{mars-idea}
lbu $t1, 0($t0)
lbu $t2, 0($s0)
xor $t1, $t1, $t2
sb $t1, 0($s1)
\end{minted}

Here \texttt{t0} points to the current keystream byte, \texttt{s0} to the source
byte, and \texttt{s1} to its destination. Increment those addresses with full-width
instructions. After exactly \(m\) iterations, subtract \(m\) from the remaining
count. If it becomes zero, return without requesting another block. Otherwise,
increment only local state word twelve and generate the next block.

A 65-byte message uses all bytes of counter \(c\), followed by the first byte of
counter \(c+1\). A 64-byte message uses only counter \(c\). Unused bytes from a final
block are discarded by this whole-message interface. Splitting the same 65-byte
message into independent calls of one and 64 bytes does not reproduce that stream
unless the caller also tracks the unused 63 bytes of the first block. A block counter
alone cannot identify a position inside a block.

For a correctness check, encrypt a byte array with the prototype and compare it
with JDK ChaCha20 using the same key, nonce and counter. Then decrypt the prototype's
result and recover the input. The round trip is useful but insufficient by itself:
two matching errors can undo each other. The independent comparison detects that
case. Include binary bytes, unaligned buffers and exact in-place buffers; printable
text alone does not cover these storage conditions.
